\chapter{Разработка концепции сервиса коллективных музыкальных рекомендаций для заведений и мероприятий}
\label{ch:service-concept}

\section{Концепция сервиса коллективных музыкальных рекомендаций для заведений и мероприятий}
\label{sec:service:concept}

В рамках настоящей работы предлагается концепция сервиса коллективных музыкальных рекомендаций для заведений и мероприятий на базе «Яндекс Музыки». Рабочее название сервиса~--- «Резонанс» или «Яндекс Музыка Резонанс». Его назначение состоит в том, чтобы формировать и динамически обновлять музыкальный поток публичного воспроизведения с учетом предпочтений посетителей, фактически находящихся в пространстве заведения или мероприятия, при сохранении управляемости со стороны бизнеса.

Базовая идея сервиса продолжает пользовательскую логику персональной «Моей Волны», но переносит ее с индивидуального слушателя на группу людей, совместно находящихся в одном физическом пространстве. Если у отдельных пользователей «Яндекс Музыки» уже есть история прослушивания, лайки, дизлайки и персональные рекомендации, то при корректном определении состава аудитории можно сформировать групповой профиль предпочтений и использовать его для генерации списка треков-кандидатов. При этом сервис не должен превращать заведение в пространство произвольного заказа музыки: бизнес сохраняет рамки по жанрам, энергичности, настроению, времени суток, стоп-листам и допустимому уровню популярности.

Для владельцев бизнеса, концертных площадок и организаторов мероприятий предполагается специальный профиль «Яндекс Музыка для бизнеса». Из этого профиля происходит настройка объекта, управление воспроизведением, выбор режимов рекомендаций, задание ограничений и получение отчетности. Музыкальный поток формируется на стороне платформы и передается в бизнес-профиль как динамический плейлист или очередь воспроизведения. Владелец заведения получает музыкальный результат, но не получает персональные данные посетителей.

Ключевым техническим предположением концепции является фиксация присутствия пользователя в заведении. Основной сценарий опирается на BLE-beacon: в пространстве размещается Bluetooth-маяк, передающий идентификатор ближайшим устройствам. Если смартфон пользователя с установленным приложением «Яндекс Музыки» находится в зоне действия маяка, приложение может отправить на сервер сигнал о нахождении пользователя в соответствующем заведении. В проектных материалах этот механизм сопоставляется с уже используемыми в экосистеме Яндекса сценариями фиксации присутствия курьера рядом с рестораном или даркстором. В качестве дополнительного или стартового сценария может использоваться QR-код: посетитель сканирует код на площадке и явно присоединяется к музыкальному пространству.

Архитектурно сервис можно описать как последовательность операций. Сначала бизнес создает профиль объекта и задает ограничения. Затем пользователь, оказавшись в заведении, получает запрос на участие и дает информированное согласие. После этого платформа формирует текущую группу пользователей, извлекает их музыкальные представления из внутреннего контура «Яндекс Музыки», строит групповой список кандидатов и применяет бизнес-ограничения. Результат передается в плеер бизнес-профиля и обновляется по мере изменения состава аудитории.

Такая концепция сознательно отделяет уровень рекомендаций от уровня промышленного лицензирования контента. Как показано в разделе~\ref{sec:market:legal}, для публичного воспроизведения музыки необходима отдельная правовая модель, включающая договоры с РАО и ВОИС и B2B-урегулирование использования каталога. Поэтому «Резонанс» следует рассматривать не просто как рекомендательный алгоритм, а как продуктовую систему, в которой рекомендательный слой, правовой слой, отчетность и пользовательское согласие должны быть согласованы.

\section{Целевые аудитории сервиса и ценностное предложение}
\label{sec:service:audiences}

У сервиса есть четыре ключевые аудитории: владельцы и управляющие заведений, организаторы мероприятий, посетители и сама экосистема Яндекса. Для каждой аудитории ценность формулируется по-разному. Бизнесу важны управляемость музыкальной атмосферы, снижение операционной нагрузки и правовая определенность. Организатору мероприятия важен более точный «разогрев» аудитории и дополнительный сигнал о ее музыкальных предпочтениях. Посетителю важна возможность мягко влиять на музыкальный фон без необходимости специально заказывать треки или раскрывать свои данные заведению. Для Яндекса сервис может усилить ценность «Яндекс Музыки», создать новый B2B-сценарий и расширить понимание контекста музыкального потребления.

\begin{table}[htbp]
\caption{Ценностное предложение сервиса для ключевых аудиторий}
\label{tab:value-proposition}
\centering
\small
\singlespacing
\begin{tabularx}{\textwidth}{p{0.22\textwidth}X X}
\toprule
Аудитория & Проблема / потребность & Ценность сервиса \\
\midrule
Владелец заведения & Музыка зависит от вкуса персонала, плейлисты устаревают, правовая модель сложна & Управляемый музыкальный поток, ограничения по стилю, автоматизированная отчетность, более точное соответствие текущей аудитории \\
\addlinespace
Организатор мероприятия & Необходимо быстро вовлечь аудиторию и подобрать уместный стартовый музыкальный материал & Групповой профиль гостей, поддержка opening-сценария, возможность адаптировать музыку без ручного опроса аудитории \\
\addlinespace
Посетитель & Хочется влиять на атмосферу, но без навязчивого заказа треков и раскрытия личных данных площадке & Добровольное участие, влияние на общую «волну», сохранение персональных данных внутри контура Яндекса \\
\addlinespace
Яндекс & Требуется расширять экосистемные сценарии и ценность музыкального сервиса & Новый B2B-продукт, повышение вовлеченности, дополнительный контекст прослушивания, развитие технологий рекомендаций \\
\bottomrule
\end{tabularx}
\end{table}

Наиболее перспективными первыми сегментами выглядят бары, рестораны, кофейни, casual dining, event-friendly пространства, клубы и организаторы специальных мероприятий. В этих форматах музыка достаточно заметна для пользовательского опыта, но при этом не всегда требует полноценного живого DJ. Менее очевидными первыми сегментами являются медицинские учреждения, аптеки, часть retail-форматов и пространства с очень строгой брендовой музыкальной айдентикой. Для них стабильный контролируемый фон может быть важнее динамической персонализации.

Ценность сервиса для бизнеса не следует формулировать как прямое обещание роста выручки. Корректнее говорить о повышении управляемости клиентского опыта, снижении риска неуместного музыкального оформления, поддержке лояльности и более точном соответствии атмосферы аудитории. Экономический эффект может проявляться через длительность пребывания, повторные визиты, удовлетворенность и вовлеченность, но эти связи требуют отдельной эмпирической проверки.

\section{Функции сервиса для бизнеса и пользователей Яндекс Музыки}
\label{sec:service:functions}

Функциональность для бизнеса строится вокруг специального профиля «Яндекс Музыка для бизнеса». В этом профиле владелец объекта или организатор мероприятия создает карточку площадки, выбирает тип пространства, указывает режим использования музыки, подключает устройство воспроизведения и задает рамки рекомендаций. Базовые ограничения могут включать жанры, уровень энергичности, допустимый диапазон темпа, нежелательные исполнители, стоп-листы, запрет на повторения, ограничения по ненормативной лексике и сценарии по времени суток.

Отдельной функцией может стать AI-консультант для настройки музыкальных ограничений естественным языком. Вместо ручного выбора десятков параметров владелец может описать задачу: «нужна спокойная музыка для кофейни утром, без тяжелого рока и слишком популярных танцевальных треков» или «нужен мягкий разогрев для вечеринки в баре, без резких переходов и с постепенным повышением энергии». Система переводит такое описание в набор ограничений, который затем применяется к рекомендательному потоку. В рамках настоящей ВКР этот элемент рассматривается как продуктовая возможность, а не как реализуемая техническая часть эксперимента.

Для бизнеса также важна отчетность. Сервис должен фиксировать фактические воспроизведения, хранить идентификаторы треков, исполнителей, авторов и изготовителей фонограмм в объеме, необходимом для подготовки отчетов РАО и ВОИС. Автоматическая генерация отчетов снижает административную нагрузку и делает продукт более ценным в сравнении с простым рекомендательным плейлистом. При этом конкретный формат отчетных файлов должен соответствовать актуальным требованиям организаций по коллективному управлению правами~\cite{rao_report_forms,vois_report_form}.

Функциональность для посетителя должна быть максимально ненавязчивой. В обычном сценарии пользователь не открывает отдельное приложение и не выбирает треки вручную. Если приложение фиксирует нахождение пользователя в пространстве, где используется сервис, пользователь получает понятное уведомление: он может «войти в резонанс» и разрешить использовать свои музыкальные предпочтения для формирования общей волны. При отказе пользователь не участвует в рекомендациях, а его пребывание не должно влиять на музыкальный поток.

Принципиальное требование к пользовательскому сценарию~--- сохранение приватности. Владелец заведения не видит список присутствующих пользователей, их историю прослушивания, лайки, дизлайки или индивидуальные музыкальные профили. Он получает только результат в виде плейлиста, очереди воспроизведения и агрегированных операционных данных. Обработка персональных данных остается внутри контура Яндекса, а пользовательское согласие должно быть явно получено и отзываемо.

\section{Пользовательские сценарии взаимодействия с сервисом}
\label{sec:service:journeys}

Сценарий владельца заведения начинается с подключения бизнес-профиля. Пользователь со стороны бизнеса выбирает тип объекта, указывает адрес, подключает устройство воспроизведения, выбирает способ фиксации присутствия посетителей и задает начальные музыкальные рамки. На старте сервис может предложить шаблон для ресторана, бара, кофейни или мероприятия, а затем уточнить параметры: уровень энергии, допустимые жанры, нежелательные направления, язык треков, режим обновления и необходимость отчетности.

В операционном режиме владелец или администратор не должен постоянно управлять музыкой вручную. Сервис формирует очередь, обновляет ее при изменении состава аудитории и показывает базовые индикаторы: активен ли маяк, сколько пользователей участвует в формировании рекомендаций в агрегированном виде, не нарушены ли ограничения, сформирован ли отчет. При необходимости администратор может временно зафиксировать плейлист, усилить или снизить энергичность, заблокировать отдельный трек или переключиться на заранее подготовленный fallback-плейлист.

Сценарий организатора мероприятия похож, но имеет более выраженную временную структуру. До события организатор создает карточку мероприятия, задает музыкальный профиль, подключает QR-код или BLE-зону и может заранее пригласить гостей «войти в резонанс». В начале события сервис формирует мягкий opening-поток, затем постепенно адаптируется к фактическому составу гостей. В кульминационной части мероприятия система может работать как источник рекомендаций для DJ или как автоматический слой для фоновых зон, не претендуя на замену живого музыкального управления.

Сценарий посетителя должен быть простым. Пользователь приходит в заведение или на мероприятие, получает уведомление или видит QR-код, знакомится с коротким объяснением и принимает или отклоняет участие. Если он соглашается, его музыкальные предпочтения учитываются в агрегированном виде. Пользователь не выбирает конкретный трек и не получает индивидуального контроля над эфиром; его вклад становится частью коллективной волны. Такая модель снижает риск конфликтов, характерный для сервисов одиночного заказа музыки, где один активный или платежеспособный гость может слишком сильно влиять на атмосферу.

Отдельный сценарий связан с использованием технологии как помощника DJ. DJ может загрузить или сопоставить подготовленный плейлист и получить оценку его потенциального соответствия аудитории, находящейся на площадке. В более развитой версии система могла бы анализировать музыкальные файлы по аудиоэмбеддингам и сопоставлять их с групповым профилем. В настоящей работе этот сценарий рассматривается как перспективное применение технологии, но не реализуется экспериментально.

Еще одно возможное применение связано с анализом вовлеченности аудитории по зонам пространства. При наличии нескольких BLE-меток теоретически можно различать посетителей, находящихся у бара, в зале или ближе к танцполу, и адаптировать музыкальный поток с учетом пространственной структуры аудитории. Такой сценарий требует особенно аккуратной работы с приватностью и не входит в техническую реализацию ВКР, но может быть полезен для дальнейшего развития продукта.

\section{Позиционирование, наименование и коммуникационная концепция сервиса}
\label{sec:service:positioning}

Рабочее название «Резонанс» удачно отражает продуктовую логику сервиса. В физическом смысле резонанс связан с усилением колебаний при совпадении частот; в коммуникационном смысле это может быть метафорой согласования индивидуальных музыкальных предпочтений многих людей в единую коллективную волну. Полное название «Яндекс Музыка Резонанс» подчеркивает связь с уже известной пользовательской экосистемой и снижает необходимость объяснять сервис как полностью новый продукт.

Для пользователя коммуникационная формула может строиться вокруг действий «Войти в резонанс» или «Резонировать». Эти формулировки не должны звучать как техническое разрешение на обработку данных; юридически необходимое согласие должно быть оформлено отдельно и понятно. В интерфейсной коммуникации важно объяснить простую пользу: пользователь может позволить заведению учитывать его музыкальный вкус при формировании общей музыки, не раскрывая свои данные владельцу площадки.

Для бизнеса позиционирование должно быть более прагматичным. Сервис не продает «магическое повышение продаж», а предлагает управляемую музыкальную среду, которая учитывает текущую аудиторию, сохраняет рамки бренда и помогает с отчетностью. Возможная формула ценности для владельца: «музыка, которая остается вашей по стилю, но лучше совпадает с гостями, которые находятся здесь сейчас».

Для организаторов мероприятий сервис может позиционироваться как инструмент мягкого вовлечения аудитории и поддержки музыкальной драматургии. В отличие от голосований и заказов треков, он не создает борьбы за эфир, а работает с агрегированным профилем группы. Это особенно важно для событий, где музыка должна поддерживать общий сценарий и не разрушаться случайными индивидуальными запросами.

Коммуникация приватности должна быть вынесена в отдельный смысловой слой. Пользователю необходимо явно сообщить, что его история прослушивания не передается заведению, не показывается администратору и не используется для идентификации перед третьими лицами. Заведение видит только музыкальный результат и агрегированные показатели. Такая формулировка снижает тревожность и помогает отличить сервис от технологий наблюдения за посетителями.

\section{Модель монетизации, дистрибуции и экономические эффекты внедрения сервиса}
\label{sec:service:monetization}

Монетизация сервиса может строиться по нескольким направлениям. Для постоянных B2B-точек наиболее естественна подписочная модель с оплатой за объект или устройство воспроизведения. Базовый тариф может включать бизнес-профиль, готовые ограничения, динамическую коллективную волну и базовую отчетность. Расширенный тариф может включать несколько зон, расширенные настройки, AI-консультанта, интеграцию с несколькими устройствами, приоритетную поддержку и более детальные отчеты. Для сетевых клиентов может использоваться индивидуальная модель с оплатой за количество объектов и дополнительными сервисными услугами.

Для событийного рынка уместна отдельная модель: оплата за событие, за день использования или процент от бюджета музыкального сопровождения. Такой подход лучше соответствует природе мероприятий, где организатору не нужна постоянная подписка, но нужна возможность быстро подключить сервис к конкретной площадке или вечеринке. Для клубов и баров с регулярными событиями возможна гибридная модель: месячная подписка для площадки плюс расширенный событийный режим.

На уровне unit-экономики базовую выручку можно описать как сумму подписочной и событийной частей:

\begin{equation*}
    R = N_{\mathrm{points}} \cdot ARPA_{\mathrm{points}} + N_{\mathrm{events}} \cdot ARPA_{\mathrm{events}},
\end{equation*}

где $N_{\mathrm{points}}$~--- число подключенных постоянных точек, $ARPA_{\mathrm{points}}$~--- средняя месячная выручка на точку, $N_{\mathrm{events}}$~--- число подключенных событий за период, а $ARPA_{\mathrm{events}}$~--- средняя выручка на событие. Эта формула не является финальным финансовым расчетом; она фиксирует принципиальное разделение двух рынков, обоснованное в разделе~\ref{sec:market:potential}.

Дистрибуция сервиса может опираться на несколько каналов. Первый~--- экосистемный канал Яндекса: продвижение среди пользователей и бизнес-клиентов, уже знакомых с «Яндекс Музыкой», «Алисой» и устройствами Яндекс Станции. Второй~--- партнерства с B2B-поставщиками музыки, POS-интеграторами, HoReCa-сервисами и event-платформами. Третий~--- прямые продажи в сегменты, где ценность коллективной музыкальной адаптации наиболее очевидна: бары, рестораны, кофейни, клубы и организаторы мероприятий. Четвертый~--- пилотные проекты с площадками, где можно продемонстрировать влияние на вовлеченность гостей и качество музыкального опыта.

Экономические эффекты для бизнеса могут проявляться в нескольких плоскостях. Во-первых, сервис снижает трудозатраты на подбор и обновление музыки. Во-вторых, он уменьшает вероятность неуместного музыкального оформления, поскольку поток адаптируется к текущей аудитории внутри заданных рамок. В-третьих, встроенная отчетность снижает административную нагрузку и правовые риски. В-четвертых, участие посетителей может повысить эмоциональную вовлеченность и создать дополнительный повод для коммуникации с брендом заведения.

Для Яндекса сервис может усилить ценность «Яндекс Музыки» как экосистемного продукта. Пользователь получает новый сценарий, в котором его музыкальная история влияет не только на персональные рекомендации, но и на реальный опыт в физическом пространстве. Бизнес получает B2B-инструмент на базе знакомой платформы. При этом у Яндекса появляется дополнительный контекст музыкального потребления, связанный не только с индивидуальным прослушиванием, но и с ситуациями совместного нахождения людей в заведениях и на мероприятиях. Такой контекст может быть ценен для дальнейшего развития рекомендательных моделей, если обработка данных осуществляется с соблюдением правовых и этических требований.

\section{Выводы по главе 2}
\label{sec:service:conclusion}

В данной главе разработана продуктовая концепция сервиса «Резонанс» / «Яндекс Музыка Резонанс» как B2B/B2C-слоя внутри экосистемы «Яндекс Музыки». Сервис предполагает специальный бизнес-профиль для заведений и мероприятий, механизм определения присутствия посетителей через BLE или QR, добровольное участие пользователей, обработку данных внутри контура Яндекса и формирование динамического группового музыкального потока.

Ключевая продуктовая идея состоит в том, что заведение сохраняет контроль над музыкальной атмосферой, а посетители получают мягкое влияние на нее без индивидуального заказа треков и без передачи персональных данных площадке. Такое решение отличается от классических B2B-сервисов музыки для бизнеса тем, что учитывает фактический состав аудитории, и отличается от сервисов заказа музыки тем, что не дает одному гостю непропорционального влияния на эфир.

Функциональность сервиса должна включать настройку рамок воспроизведения, стоп-листы, управление энергичностью, возможный AI-консультант для задания ограничений естественным языком, автоматическую генерацию отчетности для РАО и ВОИС, а также сценарии для мероприятий и поддержки DJ. Монетизация может строиться отдельно для постоянных точек и событий, поскольку эти рынки различаются по циклу продажи, LTV и логике использования.

Сформулированная концепция задает бизнес- и продуктовый контекст для последующих технических глав. Далее работа переходит к теоретическим и технологическим основам музыкальных рекомендательных систем, подходам к групповым рекомендациям и экспериментальной проверке модели на датасете YAMBDA.

% TODO(групповая ВКР): указать распределение задач между участниками проекта в рамках главы 2.
